\documentclass{beamer}
\usepackage{xeCJK}
\usepackage[tongji,zh]{collegeBeamer}
\usepackage{fontspec}
\usepackage{booktabs}
\usepackage{array}
\usepackage{tabularx}
\usepackage{graphicx}
\usepackage{tikz}
\usetikzlibrary{arrows.meta,positioning,fit,backgrounds}

\setCJKmainfont{FandolSong-Regular}[BoldFont=FandolSong-Bold,ItalicFont=FandolKai-Regular,Mapping=fullwidth-stop,Scale=0.88]
\setCJKmonofont{FandolSong-Regular}[Scale=0.82]

\setbeamertemplate{section page}{}
\AtBeginSection[]{}
\setbeamertemplate{frametitle}{
    \vspace*{-3.5ex}
    \begin{beamercolorbox}[leftskip=2cm]{frametitle}
    \usebeamerfont{frametitle}\insertframetitle
    \end{beamercolorbox}
}
\setbeamercolor{block title}{fg=white,bg=maincolor}
\setbeamercolor{block body}{fg=darkgray,bg=maincolor!4}
\setbeamerfont{block title}{series=\bfseries,size=\small}
\setbeamerfont{block body}{size=\scriptsize}

\newcommand{\smallcap}[1]{{\small\textcolor{maincolor}{#1}}}
\newcommand{\plainpara}[1]{\par\noindent #1\par}
\newcommand{\callout}[1]{\begin{center}\fcolorbox{maincolor}{maincolor!6}{\parbox{0.92\textwidth}{#1}}\end{center}}

\title{\textbf{张家浜跨模态数据处理与验证更新}}
\author{\textbf{WaterExpert}}
\date{2026-08}

\begin{document}

\begin{frame}[t]{张家浜跨模态数据处理链路}
\footnotesize
\begin{columns}[T,onlytextwidth]
\column{0.44\textwidth}
\begin{block}{输入数据}
\begin{itemize}
    \item UAV：13 资产（7 图/6 视频），覆盖 8 日。
    \item 实地监测：7 组监督信号。
    \item 历史代理：2198 水质、58370 气象，同月中位数/IQR。
    \item 近邻辅助：三鲁河段只进训练，不进目标评估。
\end{itemize}
\end{block}

\column{0.56\textwidth}
\begin{block}{处理链路}
\begin{tikzpicture}[
  font=\scriptsize,
  node distance=4.0mm and 4.0mm,
  box/.style={draw=maincolor, rounded corners=2pt, fill=white, align=center, text width=2.0cm, minimum height=0.72cm, inner sep=2.5pt},
  arrow/.style={-{Latex[length=2mm]}, thick, draw=maincolor}
]
\node[box] (raw) {UAV 图像\\UAV 视频};
\node[box, right=of raw] (feat) {抽帧\\视觉统计\\Transformer embedding};
\node[box, right=of feat, text width=2.1cm] (fusion) {日尺度跨模态表\\训练/评估隔离};
\node[box, below=of feat, text width=2.1cm] (proxy) {历史代理特征\\同月统计};
\draw[arrow] (raw) -- (feat);
\draw[arrow] (feat) -- (fusion);
\draw[arrow] (proxy) -- (fusion);
\end{tikzpicture}

\vspace{1mm}
\begin{itemize}
    \item 视频抽代表帧；图片/视频帧统一提取视觉特征。
    \item 视觉分支：Transformer 表征 + 统计特征。
    \item \texttt{sample\_site\_role} 分目标站和辅助站。
    \item 张家浜留一验证；辅助河段只做残差训练。
\end{itemize}
\end{block}
\end{columns}

\vspace{-1mm}
\smallcap{当前范围：低空 UAV 图像/视频 + 实地监测 + 历史代理特征。未接入卫星遥感。}
\end{frame}

\begin{frame}[t]{模型对比与结论}
\footnotesize
\begin{columns}[T,onlytextwidth]
\column{0.58\textwidth}
\begin{block}{张家浜留一验证}
\centering
\begin{tabular}{lcc}
\toprule
方案 & 浊度 RMSE & 透明度 RMSE \\
\midrule
非视觉 baseline & 25.9427 & 0.095851 \\
视觉残差融合 & \textbf{25.8496} & \textbf{0.095618} \\
\bottomrule
\end{tabular}
\end{block}

\begin{block}{结果}
\begin{itemize}
    \item 浊度 RMSE 下降 0.36\%。
    \item 透明度 RMSE 下降 0.24\%。
    \item 最优方案：\texttt{cross\_modal\_auxiliary\_visual\_residual}。
    \item 改进来自历史代理增强 + 视觉残差校正。
\end{itemize}
\end{block}

\column{0.42\textwidth}
\begin{block}{落库与接入}
\begin{itemize}
    \item 后端：\texttt{/api/v1/cross-modal/zhangjiabang}
    \item 媒体：\texttt{/api/v1/cross-modal/media}
    \item 前端：\texttt{frontend/visualization.html}
    \item 脚本：\texttt{frontend/js/visualization.js}
    \item 产物：日表、媒体预览、模型前后对比表
\end{itemize}
\end{block}

\begin{block}{结论}
UAV 低空视觉信息已经进入张家浜验证链路。当前是小样本稳定正收益，不表述为生产级泛化能力。
\end{block}
\end{columns}
\end{frame}

\end{document}
